\addtocontents{toc}{\protect\Needspace{3\baselineskip}}
\originalchapter{超弦自由谱与 D 膜动力学}\label{ch:branes}
\originalsection{世界面费米子}
在自由世界面加入 $\psi^I$, 使坐标与费米子具有世界面超对称. 开弦费米子可取半整数模式 NS 扇区或整数模式 R 扇区. NS 真空截距 $a=1/2$, R 截距为 0. NS 基态本含 tachyon, 选择适当 GSO 投影后保留奇数费米产生算符的态, 最低 $b^I_{-1/2}|0\rangle$ 为八个无质量玻色自由度. R 零模构成 Clifford 代数, 投影后留下一个八维手征自旋态. 这里 $D=10$ 是一致超弦的临界维数; 本卷使用该自由谱, 不展开超鬼异常和完整相互作用的推导.

闭 Type II 弦左右各可取 NS 或 R. NS--NS 与 R--R 给玻色态, NS--R 与 R--NS 给费米态, 等级匹配使左右质量贡献相等. IIA 与 IIB 的区别涉及左右 R 手征选择, 不改变本卷计算的自由能级总简并数. Heterotic 弦将超弦与玻色世界面左右结构结合, 内部规范扇区需要额外一致性构造, 此处只说明角色并列为待系统补写.

\originalsection{拉伸弦与规范场}
两平行 D$p$ 膜相距 $L$, Dirichlet 坐标含经典项 $X^\perp=y_1+L\sigma/\pi$. 长度产生基态能 $T_sL$, 质量壳变为
\begin{equation}\label{eq:stretch}
M^2=\left(\frac L{2\pi\alpha'}\right)^2+\frac{N-a}{\alpha'}.
\end{equation}
多个膜的端点标签形成 Chan--Paton 矩阵, $N$ 个重合膜的低能开弦通常给 $U(N)$ 规范场. 平行方向振子对应世界体矢量, 垂直方向对应标量. 将膜移开相当于给标量矩阵对角真空期望值, 非对角开弦获得与间距成正比的质量, 这是几何实现的 Higgs 机制.

D 膜不只是外加刚性墙. 量子谱中的无质量标量描述膜形状与位置的集体运动. 慢变化下 DBI 作用量为 $S=-T_p\int\dd^{p+1}x\sqrt{-\det(g+2\pi\alpha'F)}$, 平直背景展开给膜张力、标量动能和 Yang--Mills 动能. 此处使用低导数有效作用量, 不宣称已推导所有弦散射振幅或非 Abel 非线性 DBI 完整形式.

\originalsection{耦合与近似}
dilaton 的常值决定 $g_s=\ee^{\Phi_0}$. 世界面拓扑项使连通闭曲面亏格 $h$ 的权重为 $g_s^{2h-2}$, 盘面权重为 $g_s^{-1}$. 因此 $T_p\propto1/g_s$, 而 $G_{10}\propto g_s^2\alpha'^4$. 单膜自引力随 $g_s$ 小而弱, 堆叠 $N$ 个膜时控制参数可变为 $g_sN$. 「膜很重」与「弱耦合可忽略其自引力」并不矛盾, 因为 Newton 常数也变小.
